\section{结语}

本文以一个具体应用为目标讨论了
场空间中的平凡化映射与流。
然而，其底层概念相当普遍，
例如在严格的构造性工作、数值微扰论、
或与可重整化的光滑化技术相结合等方面，
可能还有其他用途。

所提议的威尔逊流与 HMC 算法的组合，
是否确实比迄今使用的模拟算法
更高效地对格点 QCD 的拓扑扇区采样，仍有待确定。
当格距越取越小时，
夸克场最终可能不得不被纳入流中，
或许还要包括第~4 节讨论的次领头阶修正。
无论如何，存在物质场时的平凡化映射
本身就是一个值得研究的课题。
